\documentclass[12pt,a4paper]{article}
\usepackage{goyabase}

% --- DATOS DEL DOCUMENTO ---
\title{\textbf{Guía de Conexión: DataGrip y Ecosistemas NoSQL en GloboTour}}
\author{Departamento de Informática}
\date{Abril 2026}

\usepackage[spanish, es-noshorthands, es-tabla]{babel}
\usepackage[autostyle]{csquotes}
\begin{document}

\maketitle

\section{¿Qué es Docker y para qué sirve?}
Docker es una plataforma que permite empaquetar, distribuir y ejecutar aplicaciones dentro de \textbf{contenedores}. Un contenedor es una unidad estándar de software que incluye todo lo necesario para que una aplicación funcione: código, dependencias, librerías y configuración.

\subsection{Ventajas de usar Docker en Bases de Datos}
Para un administrador de bases de datos o desarrollador, Docker ofrece beneficios críticos:
\begin{itemize}
    \item \textbf{Aislamiento:} Puedes ejecutar varios SGBD (como MongoDB y Redis) en la misma máquina sin que sus dependencias entren en conflicto.
    \item \textbf{Portabilidad:} El entorno es idéntico en tu ordenador personal, en el aula o en un servidor de producción. Si funciona en el contenedor, funciona en cualquier parte.
    \item \textbf{Limpieza:} No necesitas instalar los servicios directamente en tu sistema operativo. Cuando terminas la práctica, borras el contenedor y tu equipo queda limpio.
    \item \textbf{Rapidez de despliegue:} Con un solo comando puedes levantar una arquitectura compleja que, de forma manual, tardaría horas en configurarse.
\end{itemize}

\section*{Introducción al Entorno GloboTour}
DataGrip es un entorno de desarrollo integrado (IDE) de JetBrains diseñado específicamente para la gestión de bases de datos. A diferencia de otras herramientas que son específicas para un motor, DataGrip permite gestionar múltiples motores desde una única interfaz.

En esta práctica, simularemos la arquitectura de datos de la plataforma \textbf{GloboTour}. Conectaremos DataGrip a los servicios \textbf{MongoDB}, \textbf{Redis} y \textbf{Cassandra} desplegados mediante Docker, analizando el rol que cumple cada uno en un proyecto de alta concurrencia.

\begin{tcolorbox}[title=Requisito Previo: Gestión de Recursos]
Asegúrate de tener los contenedores en ejecución. Si tu equipo tiene **menos de 16 GB de RAM**, no levantes todos los servicios a la vez. Levanta solo el que vayas a usar y apágalo al terminar:
\begin{lstlisting}[language=bash]
# Para empezar a trabajar con uno
docker compose up -d mongodb
# Para liberar memoria antes de pasar al siguiente
docker compose stop mongodb
\end{lstlisting}
\end{tcolorbox}

\section{Conexión a MongoDB (El Catálogo Flexible)}
En GloboTour, MongoDB se utiliza para almacenar el \textbf{Catálogo Polimórfico de Experiencias}, dado que una cata de vinos y un salto en paracaídas requieren campos muy distintos.

\subsection{Alternativa: Mongo Express}
Si prefieres una interfaz web ligera, puedes acceder a través de tu navegador en \url{http://localhost:8081} (User: \texttt{admin} / Pass: \texttt{pass}).

\subsection{Pasos para la configuración en DataGrip}
\begin{enumerate}
    \item Abre DataGrip y pulsa en el icono \textbf{+} (New) $\rightarrow$ \textbf{Data Source} $\rightarrow$ \textbf{MongoDB}.
    \item En la pestaña \textbf{General}, introduce los siguientes datos:
    \begin{itemize}
        \item \textbf{Host:} \texttt{localhost}
        \item \textbf{Port:} \texttt{27017}
        \item \textbf{User/Password:} Déjalos vacíos.
    \end{itemize}
    \item Descarga los \textit{driver files} si el IDE te lo solicita en la parte inferior.
    \item Pulsa en \textbf{Test Connection} y, si aparece el check verde, haz clic en \textbf{OK}.
    \item En el panel izquierdo de DataGrip, busca la base de datos \texttt{globotour} y expande la colección \texttt{experiencias} para explorar los documentos JSON.
\end{enumerate}

\section{Conexión a Redis (Aceleración y Tiempo Real)}
Redis es nuestra capa de aceleración en memoria. Lo usaremos para gestionar la \textbf{caducidad de sesiones} de los usuarios, los \textbf{contadores de métricas} en tiempo real y los \textbf{rankings de popularidad}.

\subsection{Alternativa: RedisInsight}
Puedes usar la potente interfaz oficial en \url{http://localhost:5540}. Recuerda que para añadir la base de datos desde el navegador, el host debe ser \texttt{redis-server}.

\subsection{Pasos para la configuración en DataGrip}
\begin{enumerate}
    \item Pulsa en \textbf{+} $\rightarrow$ \textbf{Data Source} $\rightarrow$ \textbf{Redis}.
    \item Configuración de red:
    \begin{itemize}
        \item \textbf{Host:} \texttt{localhost}
        \item \textbf{Port:} \texttt{6379}
    \end{itemize}
    \item Descarga los drivers si se solicita y pulsa en \textbf{OK}.
    \item \textbf{Nota Visual:} Redis organiza los datos por bases de datos numéricas. Trabajarás principalmente en la \textbf{DB 0}.
\end{enumerate}

\begin{tcolorbox}[title=Consejo de Productividad: Jerarquía Visual]
En la práctica de GloboTour, usamos el carácter \enquote{:} en las llaves (ej: \texttt{session:user:101} o \texttt{experiencias:top\_ventas}). DataGrip interpretará estos dos puntos como carpetas, generando una estructura jerárquica que te facilitará la navegación visual.
\end{tcolorbox}

\section{Conexión a Apache Cassandra (Registro de Búsquedas)}
Usamos Cassandra para almacenar \textbf{millones de registros de búsqueda diarios}. Al ser de tipo columnar y altamente distribuida, permite escrituras masivas sin bloquear la base de datos principal, siguiendo un diseño orientado a las consultas (\textit{Query-First Design}).

\subsection{Alternativa: CloudBeaver}
Puedes usar el gestor profesional CloudBeaver en \url{http://localhost:8978}. Al conectar, usa el host \texttt{cassandra-server}.

\subsection{Pasos para la configuración en DataGrip}
\begin{enumerate}
    \item Pulsa en \textbf{+} $\rightarrow$ \textbf{Data Source} $\rightarrow$ \textbf{Cassandra}.
    \item Configuración:
    \begin{itemize}
        \item \textbf{Host:} \texttt{localhost}
        \item \textbf{Port:} \texttt{9042}
    \end{itemize}
    \item Asegúrate de descargar o actualizar los drivers para soportar correctamente las sentencias CQL.
    \item Pulsa en \textbf{Test Connection}. Ten en cuenta que Cassandra suele tardar un poco más en arrancar que los otros contenedores; si falla a la primera, espera unos segundos e inténtalo de nuevo.
    \item Aquí gestionarás, por ejemplo, la creación de la tabla \texttt{busquedas\_por\_fecha} diseñando adecuadamente las \textit{Partition Keys}.
\end{enumerate}

\section{Conexión a Neo4j (La Red de Recomendaciones)}
Neo4j es nuestra base de datos de \textbf{Grafos}. La usamos para encontrar relaciones complejas (amigos de amigos, rutas sugeridas) que en SQL requerirían procesos muy lentos.

\subsection{Alternativa: Neo4j Browser}
Neo4j incluye una consola interactiva muy potente para visualizar grafos en \url{http://localhost:7474}.

\subsection{Pasos para la configuración en DataGrip}
\begin{enumerate}
    \item Pulsa en \textbf{+} $\rightarrow$ \textbf{Data Source} $\rightarrow$ \textbf{Neo4j}.
    \item Configuración:
    \begin{itemize}
        \item \textbf{Host:} \texttt{localhost}
        \item \textbf{Port:} \texttt{7687} (Protocolo Bolt)
        \item \textbf{User/Password:} \texttt{neo4j} / \texttt{password123}
    \end{itemize}
    \item Pulsa en \textbf{Test Connection}. DataGrip te permitirá escribir consultas en lenguaje \textbf{Cypher} y ver el resultado de forma gráfica o tabular.
\end{enumerate}

\section{Ventajas de usar DataGrip en NoSQL}
\begin{itemize}
    \item \textbf{Consola unificada:} Te permite escribir sentencias CQL, comandos de Redis o pipelines de agregación de Mongo desde una sola herramienta.
    \item \textbf{Exploración visual:} Puedes editar BSON/JSON directamente en la rejilla sin escribir tediosos comandos de actualización.
    \item \textbf{Exportación rápida:} Facilita la exportación de los datos analizados a CSV o JSON para incluirlos en el informe final de la práctica.
\end{itemize}

\end{document}
